\documentclass[11pt]{article}
\usepackage[margin=1in]{geometry}
\usepackage{amsmath,amssymb,amsthm,mathtools}
\usepackage{microtype}
\usepackage[hidelinks]{hyperref}

\newtheorem{theorem}{Theorem}[section]
\newtheorem{lemma}[theorem]{Lemma}
\newtheorem{corollary}[theorem]{Corollary}
\newtheorem{remark}[theorem]{Remark}
\newcommand{\C}{\mathbb C}
\newcommand{\Tr}{\operatorname{Tr}}
\newcommand{\ran}{\operatorname{ran}}
\newcommand{\rank}{\operatorname{rank}}
\newcommand{\id}{I}
\newcommand{\join}{\mathbin{\vee}}

\title{The unrestricted quantum chromatic number of the Lalonde graph joins}
\author{Independent research draft}
\date{August 1, 2026}

\begin{document}
\maketitle

\begin{abstract}
Let $G_{19}$ be the 19-vertex complex orthogonality graph introduced by
Lalonde and put $G_n=G_{19}\join K_{n-3}$ for $n\geq3$.  It was known that
$\chi_q^{(1)}(G_n)=n+1$ and conjectured that the same is true without a rank
restriction.  We prove the conjecture in the finite-dimensional projector
model:
\[
   \chi_q(G_n)=n+1\qquad(n\geq3).
\]
In particular, adjoining one apex to $G_{19}$ gives a 20-vertex graph $H$
with $\chi_q(H)=5$.

The proof first symmetrizes an arbitrary coloring, so that every projector
has a common rank $r$.  For a fixed color, a rational sum-of-squares identity
rigidifies vertices $1,\ldots,13$ into a rank-$r$ sign-vector configuration.
The remaining six subspaces are classified, including all non-transverse
branches, by an $r$-plane invariant under a canonical complex structure.
Orthogonality between two colors then preserves the two complex-structure
sectors while reversing membership in the invariant planes.  The resulting
sector subspaces are pairwise orthogonal across colors, and their dimensions
give the impossible inequality $3nr\leq2nr$.  The core sum-of-squares and all
formal block identities have an exact, standard-library verifier.
\end{abstract}

\section{Graph, model, and result}

The vertex set of $G_{19}$ is $\{1,\ldots,19\}$, and its edge set is
\begin{align*}
E_{19}=\{&12,13,18,19,1\,15,
23,26,27,2\,14,34,35,3\,16,45,4\,12,4\,13,\\
&5\,10,5\,11,5\,17,67,6\,11,6\,13,
7\,10,7\,12,7\,19,89,8\,11,8\,12,8\,18,\\
&9\,10,9\,13,14\,17,14\,18,15\,17,15\,19,16\,18,16\,19\}.
\end{align*}
Here $vw$ denotes the unordered edge $(v,w)$.  Its graph6 representation is
\begin{center}
\texttt{RxLAKA@AgYAWDGO?O?@??A?W@@OC@\_}.
\end{center}
The only triangles are
\[
 T_0=123,\qquad T_1=189,\qquad T_2=345,\qquad T_3=267.
\]
For $n\geq3$, define
\[
   G_n=G_{19}\join K_{n-3}.
\]
Thus the graph in the question is $H=G_4$, with vertex 20 as the joined
apex.

An $m$-coloring in the finite-dimensional projector model is a family of
orthogonal projections $P_{v,c}\in M_d(\C)$, with zero projections allowed,
such that
\[
 \sum_{c=1}^mP_{v,c}=\id_d,
 \qquad
 P_{v,c}P_{w,c}=0\quad(vw\in E).
\]
The row orthogonality $P_{v,c}P_{v,c'}=0$ for $c\ne c'$ follows from, and may
equivalently be included with, the assertion that the row is a projective
partition of the identity.

\begin{theorem}[Main theorem]\label{thm:main}
For every $n\geq3$,
\[
       \chi_q(G_{19}\join K_{n-3})=n+1.
\]
In particular, $\chi_q(H)=5$.
\end{theorem}

For the elementary upper bound, $G_{19}$ has the following four-coloring:
\[
\begin{array}{c|rrrrrrrrrrrrrrrrrrr}
v&1&2&3&4&5&6&7&8&9&10&11&12&13&14&15&16&17&18&19\\ \hline
\kappa(v)&1&2&3&2&1&3&1&2&3&2&4&3&1&1&2&1&3&3&3.
\end{array}
\]
Giving the $n-3$ joined vertices fresh colors proves $\chi(G_n)\leq n+1$.
The work is to exclude a quantum $n$-coloring.  This also excludes every
coloring with fewer than $n$ colors, since unused colors may be padded by
zero projections.

\section{Clique saturation and uniformization}

\begin{lemma}[Clique columns]\label{lem:clique}
Suppose an $m$-coloring contains an $m$-clique $C$.  Then, for every color
$c$,
\[
        \sum_{v\in C}P_{v,c}=\id_d.
\]
Thus the $m\times m$ array $(P_{v,c})_{v\in C,c\in[m]}$ is a magic unitary.
\end{lemma}

\begin{proof}
For fixed $c$, the summands are pairwise orthogonal, so their sum is a
projection of trace at most $d$.  Summing this trace bound over the $m$
columns gives at most $md$.  The row identities show that the total trace is
exactly $md$, so equality holds in every column.  A projection of trace $d$
in $M_d(\C)$ is the identity.
\end{proof}

\begin{lemma}[Color uniformization]\label{lem:uniform}
If $G_n$ has a finite-dimensional quantum $n$-coloring, then it has one in a
dimension $D=nr$ for which every projector has the same positive rank $r$.
\end{lemma}

\begin{proof}
Given a coloring in dimension $d$, take its direct sum over the $n!$ color
permutations:
\[
 \widetilde P_{v,c}=\bigoplus_{\pi\in S_n}P_{v,\pi(c)}.
\]
This is again a coloring.  Its dimension is $D=n!d$, while
\[
 \rank \widetilde P_{v,c}
  =(n-1)!\sum_{a=1}^n\rank P_{v,a}=(n-1)!d=:r.
\]
Hence $D=nr$ and $r>0$.
\end{proof}

Assume henceforth that such a uniform coloring exists.  Fix a color $c$.
Let $R_c$ be the sum of the color-$c$ projections on the joined clique
$K_{n-3}$ and let
\[
    Q_c=\id_D-R_c.
\]
The joined vertices are pairwise adjacent, so $\rank R_c=(n-3)r$ and
$\dim Q_c=3r$.  Every old-vertex projection $P_{v,c}$ is supported on $Q_c$.
Applying Lemma~\ref{lem:clique} to $T_j\join K_{n-3}$ gives
\begin{equation}\label{eq:triangles}
       \sum_{v\in T_j}P_{v,c}=Q_c
       \qquad(j=0,1,2,3).
\end{equation}
Thus, on $Q_c$, the fixed-color projections form a uniform rank-$r$
projective orthogonal representation of $G_{19}$ in dimension $3r$.

\section{Exact fixed-color rigidity}

In this section fix the color and write $e_v=P_{v,c}|_{Q_c}$ and
$I=I_{Q_c}$.  All traces are the ordinary matrix trace on $Q_c$.

\subsection{A rational core sum of squares}

Put
\[
 X=e_{10}+e_{11}+e_{12}+e_{13}
\]
and define the three Walsh forms
\begin{align*}
 H_1&=e_{10}+e_{11}-e_{12}-e_{13},\\
 H_2&=e_{10}-e_{11}+e_{12}-e_{13},\\
 H_3&=e_{10}-e_{11}-e_{12}+e_{13}.
\end{align*}
Also put $d_1=e_5-e_4$, $d_2=e_7-e_6$, and $d_3=e_9-e_8$.

\begin{lemma}[Core SOS]\label{lem:sos}
The defining projection, edge, and triangle relations imply the exact
identity
\begin{equation}\label{eq:sos}
       \frac43I-X=\sum_{j=0}^3 f_j^*f_j,
\end{equation}
where all $f_j$ are self-adjoint and
\[
 f_0=\frac12\left(X-\frac43I\right),
 \qquad
 f_i=\frac23d_i+\frac12H_i\quad(1\leq i\leq3).
\]
\end{lemma}

\begin{proof}
The triangle decompositions and edge zeros give
\begin{equation}\label{eq:pair-squares}
 \sum_{i=1}^3d_i^2
  =(e_4+e_5)+(e_6+e_7)+(e_8+e_9)=2I.
\end{equation}
Hadamard cancellation, together with $e_j^2=e_j$, gives
\begin{equation}\label{eq:walsh-squares}
       X^2+H_1^2+H_2^2+H_3^2=4X.
\end{equation}
Finally,
\begin{equation}\label{eq:neighbor-anti}
       \sum_{i=1}^3\{d_i,H_i\}=-4X.
\end{equation}
To see the last identity, collect the coefficient of each $e_j$ for
$j=10,11,12,13$.  The corresponding signed sum of $d_1,d_2,d_3$ is
$N_j-(2I-N_j)=2N_j-2I$, where $N_j$ is the sum of the three projections among
$e_4,\ldots,e_9$ adjacent to $j$.  Since $N_je_j=e_jN_j=0$, its
anticommutator with $e_j$ is $-4e_j$.  Summing over $j$ proves
\eqref{eq:neighbor-anti}.

Expanding the right side of \eqref{eq:sos} and inserting
\eqref{eq:pair-squares}--\eqref{eq:neighbor-anti} leaves exactly
$\frac43I-X$.
\end{proof}

The two sides of \eqref{eq:sos} have the same trace, namely zero, because
$\dim Q_c=3r$ and each $e_j$ has rank $r$.  Faithfulness of the matrix trace
therefore forces every $f_j$ to vanish.  Define
\[
       A=e_4-e_5,\qquad B=e_6-e_7,\qquad C=e_8-e_9.
\]
Walsh inversion now gives
\begin{align}
e_{10}&=\tfrac13(I+A+B+C),&
e_{11}&=\tfrac13(I+A-B-C),\label{eq:walsh-inverse}\\
e_{12}&=\tfrac13(I-A+B-C),&
e_{13}&=\tfrac13(I-A-B+C).\nonumber
\end{align}
Moreover,
\[
 A^2=I-e_3,\qquad B^2=I-e_2,\qquad C^2=I-e_1,\qquad
 A^2+B^2+C^2=2I.
\]
Inserting \eqref{eq:walsh-inverse} into $e_j^2=e_j$ and separating the four
Walsh sign patterns yields
\begin{equation}\label{eq:anticommutators}
        \{B,C\}=A,\qquad \{A,C\}=B,\qquad \{A,B\}=C.
\end{equation}

\begin{lemma}[Core normal form]\label{lem:core}
There is a Hilbert space $K_c$ of dimension $r$ and a unitary
$W_c:K_c^3\to Q_c$ under which the ranges of $e_1,\ldots,e_{13}$ are
\begin{align*}
1&=(I,0,0),&2&=(0,I,0),&3&=(0,0,I),\\
4&=(I,I,0),&5&=(I,-I,0),\\
6&=(I,0,I),&7&=(I,0,-I),\\
8&=(0,I,I),&9&=(0,I,-I),\\
10&=(I,I,I),&11&=(I,I,-I),\\
12&=(I,-I,I),&13&=(I,-I,-I).
\end{align*}
Here a block row denotes the range of the associated block-column map from
$K_c$; normalization of the frame is immaterial to its range.
\end{lemma}

\begin{proof}
Use $e_1+e_2+e_3=I$ to write $Q_c=K_1\oplus K_2\oplus K_3$, with
$K_i=\ran e_i$ and $\dim K_i=r$.  The operator $A$ is supported on
$K_1\oplus K_2$, $B$ on $K_1\oplus K_3$, and $C$ on
$K_2\oplus K_3$.  In the first relation of \eqref{eq:anticommutators}, the
$(1,1)$ and $(2,2)$ blocks of $\{B,C\}$ vanish, so the two diagonal blocks of
$A$ vanish.  Cyclically, all diagonal blocks vanish.  Since
$A^2=I-e_3$, its $(1,2)$ block is unitary, and similarly for $B,C$.  Thus
\[
A=\begin{pmatrix}0&U&0\\U^*&0&0\\0&0&0\end{pmatrix},\quad
B=\begin{pmatrix}0&0&V\\0&0&0\\V^*&0&0\end{pmatrix},\quad
C=\begin{pmatrix}0&0&0\\0&0&W\\0&W^*&0\end{pmatrix}.
\]
Equation \eqref{eq:anticommutators} gives
$U=VW^*$, $V=UW$, and $W=U^*V$.  Conjugating by the block-diagonal unitary
$\operatorname{diag}(I,U^*,V^*)$ makes all three off-diagonal blocks the
identity.  The binary decompositions and \eqref{eq:walsh-inverse} then give
the displayed ranges.
\end{proof}

\subsection{All tail branches}

The tail is not always the graph of an operator; retaining degenerate planes
is essential.

\begin{lemma}[Tail-plane classification]\label{lem:tail}
After the normalization in Lemma~\ref{lem:core}, the ranges at vertices
$14,\ldots,19$ are determined by an $r$-plane
$M_c\subset K_c\oplus K_c$ satisfying
\begin{equation}\label{eq:J-invariance}
       J_cM_c=M_c,\qquad
       J_c=\begin{pmatrix}0&I\\-I&0\end{pmatrix}.
\end{equation}
Explicitly,
\begin{align}
14&=\{(a,0,b):(a,b)\in M_c\},&
15&=\{(0,a,b):(a,b)\in M_c\},\nonumber\\
16&=\{(a,-b,0):(a,b)\in M_c\},\label{eq:tail-form}\\
17&=\{(x,x,z):(x,z)\in M_c^\perp\},&
18&=\{(x,-z,z):(x,z)\in M_c^\perp\},\nonumber\\
19&=\{(x,-z,x):(x,z)\in M_c^\perp\}.\nonumber
\end{align}
Conversely, every $J_c$-invariant $r$-plane gives rank-$r$ subspaces
satisfying all fixed-color edge relations.
\end{lemma}

\begin{proof}
The six core--tail edges $2\,14$, $1\,15$, $3\,16$, $5\,17$, $8\,18$, and
$7\,19$ place the tail ranges in the images of the following injective
two-column frames:
\[
\begin{array}{c|c@{\qquad}c|c}
v&F_v(a,b)&v&F_v(a,b)\\ \hline
14&(a,0,b)&17&(a,a,b)\\
15&(0,a,b)&18&(a,-b,b)\\
16&(a,-b,0)&19&(a,-b,a).
\end{array}
\]
Write $N_v\subset K_c^2$ for the $r$-dimensional parameter plane at $v$.
For the six tail--tail edges, direct pullback gives
\[
\begin{array}{c|cccccc}
vw&14\,17&15\,17&14\,18&16\,18&16\,19&15\,19\\ \hline
F_v^*F_w&I&I&I&I&I&-J_c.
\end{array}
\]
Put $M_c=N_{14}$.  The first five entries successively give
$N_{17}=M_c^\perp$, $N_{15}=M_c$, $N_{18}=M_c^\perp$,
$N_{16}=M_c$, and $N_{19}=M_c^\perp$.  The final entry says
$J_cM_c^\perp=M_c^\perp$.  Because $J_c$ is unitary and $J_c^*=-J_c$,
this is equivalent to $J_cM_c=M_c$.  The displayed frames are exactly
\eqref{eq:tail-form}.  Reversing these equalities verifies the converse.
\end{proof}

\begin{remark}
If the first-coordinate projection of $M_c$ is invertible, then
$M_c=\ran(I,S)$ and \eqref{eq:J-invariance} is $S^2=-I$.  This transverse
branch recovers the familiar operator-valued normal form.  It is not
exhaustive.  For instance,
$M=\operatorname{span}\{(e_1,0),(0,e_1)\}\subset\C^2\oplus\C^2$ is
$J$-invariant but is not a graph.  Lemma~\ref{lem:tail} includes this branch.
\end{remark}

\section{Cross-color sector obstruction}

Take distinct colors $c,d$ and use the core isometries from
Lemma~\ref{lem:core}.  Put
\[
       T_{dc}=W_d^*W_c:K_c^3\longrightarrow K_d^3.
\]

\begin{lemma}[Cross-color block equations]\label{lem:cross}
There are maps $X,Y,Z:K_c\to K_d$ such that
\begin{equation}\label{eq:skew-block}
T_{dc}=\begin{pmatrix}
0&X&Y\\-X&0&Z\\-Y&-Z&0
\end{pmatrix}.
\end{equation}
For $L:K_c\to K_d$, define
\[
        \Omega_L=\begin{pmatrix}0&L\\-L&0\end{pmatrix}.
\]
Then, for each $L\in\{X,Y,Z\}$,
\begin{equation}\label{eq:plane-flip}
 \Omega_LM_c\subseteq M_d^\perp,
 \qquad
 \Omega_LM_c^\perp\subseteq M_d.
\end{equation}
\end{lemma}

\begin{proof}
Same-vertex orthogonality at vertices $1,2,3$ kills the diagonal blocks of
$T_{dc}$.  Vertices $4,6,8$ make the corresponding lower block the negative
of the upper block, giving \eqref{eq:skew-block}.  No adjoints occur here:
$T_{dc}$ maps the coordinate space for $c$ to the one for $d$.

Compressing \eqref{eq:skew-block} by the two-column frames in
\eqref{eq:tail-form}, vertices $14,15,16$ give, respectively,
\[
 \Omega_YM_c\subseteq M_d^\perp,\qquad
 \Omega_ZM_c\subseteq M_d^\perp,\qquad
 \Omega_XM_c\subseteq M_d^\perp.
\]
Vertices $17,18,19$ give
\[
 \Omega_{Y+Z}M_c^\perp\subseteq M_d,\qquad
 \Omega_{Y-X}M_c^\perp\subseteq M_d,\qquad
 \Omega_{Z-X}M_c^\perp\subseteq M_d.
\]
The coefficient transform is invertible; for example
\[
 2X=(Y+Z)-(Y-X)-(Z-X),
\]
with analogous formulas for $Y,Z$.  This proves \eqref{eq:plane-flip}.
\end{proof}

Let $\Gamma_c=2P_{M_c}-I$ on $K_c^2$.  Equation
\eqref{eq:plane-flip} is equivalent to
\begin{equation}\label{eq:gamma-anti}
      \Gamma_d\Omega_L=-\Omega_L\Gamma_c
      \qquad(L=X,Y,Z).
\end{equation}
Because $M_c$ reduces the unitary $J_c$, it splits across the $+i$ and $-i$
eigenspaces of $J_c$.  Using the isometries up to a harmless factor
\[
       V_{c,\sigma}(x)=(x,\sigma i x),\qquad \sigma\in\{+,-\},
\]
write
\[
 M_c\cap\ker(J_c-\sigma iI)=V_{c,\sigma}(E_{c,\sigma}),
 \qquad e_{c,\sigma}=\dim E_{c,\sigma}.
\]
Then
\begin{equation}\label{eq:sector-rank}
          e_{c,+}+e_{c,-}=r.
\end{equation}
The operator $\Omega_L$ intertwines $J_c$ with $J_d$ and acts within a
$\sigma$ sector as the nonzero scalar $\sigma i$ times $L$.  Combining this
with \eqref{eq:gamma-anti} gives
\begin{equation}\label{eq:sector-map}
          L E_{c,\sigma}\subseteq E_{d,\sigma}^{\perp}
          \qquad(c\ne d,\ L=X,Y,Z).
\end{equation}

We can now finish the proof.

\begin{proof}[Proof of Theorem~\ref{thm:main}]
Suppose that $G_n$ had a quantum $n$-coloring.  Uniformize it using
Lemma~\ref{lem:uniform}, so the physical dimension is $D=nr$ and all ranks
are $r>0$.  For each color $c$, apply Lemmas~\ref{lem:core} and
\ref{lem:tail}, and define the physical sector subspace
\[
       \mathcal H_{c,\sigma}
       =W_c(E_{c,\sigma}\oplus E_{c,\sigma}\oplus E_{c,\sigma}).
\]
It has dimension $3e_{c,\sigma}$.  Equations
\eqref{eq:skew-block} and \eqref{eq:sector-map} show that, for fixed
$\sigma$, the $n$ subspaces $\mathcal H_{c,\sigma}$ are pairwise orthogonal
in the common $nr$-dimensional physical space.  Therefore
\[
       3\sum_{c=1}^n e_{c,+}\leq nr,
       \qquad
       3\sum_{c=1}^n e_{c,-}\leq nr.
\]
Adding and using \eqref{eq:sector-rank} gives
\[
        3nr=3\sum_{c=1}^n(e_{c,+}+e_{c,-})\leq2nr,
\]
which is impossible because $n,r>0$.  Hence no quantum $n$-coloring exists.
The explicit classical $(n+1)$-coloring given in Section 1 proves equality.
\end{proof}

\section{Relation to the rank-one theorem and exact verification}

The previously known rank-one argument classifies scalar three-dimensional
orthogonal representations of $G_{19}$.  The proof above does not assume that
classification and does not assume equal ranks in the original coloring.
Instead, color symmetrization makes the fixed-color ranks uniform; the core
SOS classifies the resulting fusion-frame core; and
Lemma~\ref{lem:tail} treats every higher-rank tail, including non-graph and
noncommuting branches.  Thus the obstruction applies to arbitrary finite
dimension, zero projectors, nonuniform original ranks, reducible
representations, and irreducible representations alike.

The machine-readable certificate and verifier distributed with this draft
replay, using exact rational arithmetic, the free-algebra expansion of
Lemma~\ref{lem:sos}, its reduction to the graph relations, all six tail
compressions in Lemma~\ref{lem:cross}, and the final integer dimension
arithmetic.  No numerical SDP output enters the theorem.

\begin{thebibliography}{9}
\bibitem{Lalonde2026}
O.~Lalonde,
\emph{Quantum colorings of spheres},
\href{https://arxiv.org/abs/2606.10872}{arXiv:2606.10872v1} (2026).
\end{thebibliography}

\end{document}
